% Generated draft: requires mathematical and source review.
\documentclass[11pt]{article}
\usepackage[margin=25mm]{geometry}
\usepackage{fontspec}
\setmainfont{Latin Modern Roman}
\usepackage{amsmath,amssymb,mathtools,unicode-math}
\setmathfont{Latin Modern Math}
\usepackage{enumitem,needspace,xurl,hyperref}
\setlength{\parindent}{0pt}
\setlength{\parskip}{5pt}
\setlength{\emergencystretch}{3em}
\newcommand{\N}{\mathbb{N}}
\newcommand{\Z}{\mathbb{Z}}
\newcommand{\Q}{\mathbb{Q}}
\newcommand{\R}{\mathbb{R}}
\newcommand{\C}{\mathbb{C}}
\newcommand{\F}{\mathbb{F}}
\newcommand{\A}{\mathbb{A}}
\newcommand{\PP}{\mathbb{P}}
\newcommand{\E}{\mathbb{E}}
\newcommand{\eps}{\varepsilon}
\newcommand{\dd}{\,\mathrm{d}}
\begin{document}
\section*{UnsolvedMath ID: 20003096}
\textbf{Source problem code:} catalogue entry
\subsection{Definitions and mathematical statement}
K. Rafi asks whether the behavior of Thurston geodesics can be understood inductively by passing from a surface to its subsurfaces. In the usual finite-type setting, let $\mathcal T(S)$ carry the ordered asymmetric Thurston metric
\[
d_{\mathrm{Th}}(X,Y)=\log\sup_{\alpha\in\mathcal C(S)}\frac{\ell_Y(\alpha)}{\ell_X(\alpha)}.
\]
A forward unit-speed geodesic $t\mapsto X_t$ satisfies $d_{\mathrm{Th}}(X_s,X_t)=t-s$ for $s\leq t$. The source does not specify a precise induction scheme or quantitative reconstruction target.

The aid reports that the 2024 active-interval theorem of Lenzhen--Modami--Rafi--Tao gives a coarse affirmative description: relevant subsurface shadows move without backtracking during active intervals with bounded boundary length and have bounded movement outside those intervals. Exact subsurface, projection, relevance and constant hypotheses are not supplied; this is not promoted to a formula for recursively reconstructing a global geodesic from one subsurface.

Its separate stretch-defect lemma is formulated exactly. Fix an essential subsurface $R$ and a nonempty fixed set $\mathcal C(R)\subseteq\mathcal C(S)$ of essential curve classes carried by it. For $s\leq t$ set
\[
L_R(X_s,X_t)=\log\sup_{\alpha\in\mathcal C(R)}
\frac{\ell_{X_t}(\alpha)}{\ell_{X_s}(\alpha)},\qquad
\Delta_R(s,t)=(t-s)-L_R(X_s,X_t).
\]
The lengths are those of the ambient marked structures. This restricted functional need not be nonnegative and is not asserted to be an intrinsic Thurston metric on $R$. Nevertheless its defect along the given geodesic is nonnegative, and for $s\leq t\leq u$ the aid asserts
\[
\Delta_R(s,u)\geq\Delta_R(s,t)+\Delta_R(t,u)\geq0.
\]
Hence for every $\eps\geq0$,
\[
\Delta_R(s,u)\leq\eps\quad\Longrightarrow\quad
\Delta_R(p,q)\leq\eps
\quad(s\leq p\leq q\leq u).
\]
Exact full-speed visibility in $R$ is therefore hereditary on nested intervals, and the same absolute error bound survives without enlargement. This is an absolute-error statement, not a uniform relative-speed estimate on arbitrarily short subintervals.

The aid also gives an obstruction to using a single proper subsurface as complete distance data: a mapping-class construction keeps its restricted length spectrum unchanged while global distance diverges. For example in the applicable disjoint-support situation, a pseudo-Anosov supported in an essential positive-complexity subsurface disjoint from $R$ fixes every class in $\mathcal C(R)$; the orbit of a marked structure then has $L_R$ equal to zero relative to the initial structure while global displacement tends to infinity. This endpoint blindness does not assert that $R$ is inactive at every time on every joining geodesic. A reconstruction theory needs information from multiple subsurfaces and compatibility among it.

The supplied verification identifies the post-separator paragraph about dynamical Teichm\"uller geodesic flow, billiards and flat structures as context introducing Problem 2.7, not part of Problem 2.6. It is therefore not treated as an additional task or hypothesis here.
\subsection{Short English statement}
Understand Thurston geodesics through subsurface data. The supplied aid reports coarse active-interval theory and a hereditary restricted-stretch defect, while explaining why one proper subsurface alone cannot recover global distance.
\subsection{Sources}
\begin{itemize}
\item Local UnsolvedMath registry, permanent ID $20003096$, source code \texttt{catalogue entry}; canonical \texttt{aim-topology-notes.json}, \texttt{notes[183]}.
\item AIM workshop \emph{Lipschitz metric on Teichmueller space}, \emph{Problems on Thurston Metric}, Problem 2.6, attributed to K. Rafi: \url{https://aimath.org/pastworkshops/lipschitzteichproblems.pdf}. Original remarks and literature field are empty. The supplied source-verification report separates the subsequent Problem 2.7 introduction from this question.
\item Lenzhen--Modami--Rafi--Tao, 2024 active-interval theorem, as cited in the aid; no title, exact theorem number or authoritative full bibliographic entry is supplied.
\item Research attempt 1: \texttt{partial\_result}, status \texttt{partially\_solved}, classification \texttt{PARTIAL-PROGRESS}, rigor \texttt{proved}, novelty confidence low. Its defect and blindness claims are machine-generated aid, not peer reviewed, and require independent verification before citation.
\end{itemize}
\end{document}
